\chapter{最適割当問題}
\label{ch:sem-assignment-chapter}

\begin{remark}{本章の設定}{sem-ch2-setup}
  $(\X, d_\X)$, $(\Y, d_\Y)$ は\textbf{完備可分な距離空間}（Polish 空間，
  定義~\ref{def:sem-polish}）とする．
  定義~\ref{def:sem-borel} $\Bb(\X)$ により
  $\X$ を可測空間（定義~\ref{def:sem-measurable-space}）$(\X, \Bb(\X))$ として扱う
  （$\Y$ についても同様）．$\X$ 上の確率測度（定義~\ref{def:sem-probability-measure}）全体を $\Mm_+^1(\X)$ と記す．

  Polish 空間とその構成要素（距離・開球・開集合・ボレル $\sigma$-代数・完備性・
  可分性）は付録（前提知識）にまとめてある．離散版で扱う有限点集合 $S = \{x_1, \ldots, x_n\} \subset \X$ は，
  どの距離を入れても自動的に Polish 空間となる（注意~\ref{rem:sem-finite-auto-polish}）．
  したがって本章の Polish 仮定は連続側でのみ効き，離散版には追加制約を課さない．
\end{remark}

整数の集合 $\range{n} = \{1, \ldots, n\}$ 上の\textbf{置換}
（全単射 $\sigma \colon \range{n} \to \range{n}$，定義~\ref{def:sem-permutation}）の
全体を $\Perm(n)$ と書く．$\Perm(n)$ は $|\Perm(n)| = n!$ の有限集合である．

\begin{definition}{最適割当問題}{sem-assignment}
  $n \in \N$，置換 $\sigma \in \Perm(n)$，
  および行列 $\mathbf{C} \in \Rnn^{n \times n}$ に対して，
  \[
    \min_{\sigma \in \Perm(n)} \frac{1}{n} \sum_{i=1}^{n} C_{i, \sigma(i)}
  \]
  を達成する $\sigma$ を求める問題を\textbf{最適割当問題}という．
\end{definition}

添字 $i$ から $j$ への\textbf{輸送コスト}を $C_{i,j}$ と解釈すれば，
$\sigma(i) = j$ は $i$ を $j$ に割り当てることを意味し，
目的関数 $\frac{1}{n}\sum_i C_{i, \sigma(i)}$ はその割当の平均輸送コストを表す．

\begin{example}{最適割当}{sem-assignment-cost-example}
  $n = 2$ とし，置換 $\sigma \in \Perm(2)$ と行列
  \[
    \mathbf{C} =
    \begin{pmatrix}
      1 & 4 \\
      2 & 1
    \end{pmatrix} \in \Rnn^{2 \times 2}
  \]
  を考える．$\Perm(2) = \{(1,2),\, (2,1)\}$ の $2! = 2$ 通りの割当を以下に図示する．

  \begin{center}
  \begin{tikzpicture}[
    wk/.style={circle, draw=blue!60, fill=blue!10,
               minimum size=18pt, inner sep=1pt, font=\small},
    tk/.style={rectangle, rounded corners=1.5pt,
               draw=red!60, fill=red!10,
               minimum size=18pt, inner sep=1pt, font=\small},
    optarr/.style={-{Stealth[length=5pt]}, thick, green!50!black},
    arr/.style={-{Stealth[length=5pt]}, semithick, black!60, dashed},
    costlbl/.style={fill=white, inner sep=1pt, font=\footnotesize},
  ]
    \node[wk] (w1) at (0, 1.2) {1};
    \node[wk] (w2) at (3, 1.2) {2};
    \node[tk] (t1) at (1, -1.2) {1};
    \node[tk] (t2) at (4, -1.2) {2};

    % σ=(1,2): 1→1, 2→2（最適）
    \draw[optarr] (w1) -- node[costlbl, midway] {(a)\ $1$} (t1);
    \draw[optarr] (w2) -- node[costlbl, midway] {(a)\ $1$} (t2);
    % σ=(2,1): 1→2, 2→1
    \draw[arr] (w1) -- node[costlbl, pos=0.25] {(b)\ $4$} (t2);
    \draw[arr] (w2) -- node[costlbl, pos=0.25] {(b)\ $2$} (t1);
  \end{tikzpicture}
  \end{center}

  各割当のコストをまとめると，
  \begin{center}
  \begin{tabular}{ccc}
     & $\sigma$ & コスト $\frac{1}{n}\sum_i C_{i,\sigma(i)}$ \\\hline
    (a) & $(1,2)$ & $\frac{1}{2}(1 + 1) = \mathbf{1}$ \\
    (b) & $(2,1)$ & $\frac{1}{2}(4 + 2) = 3$ \\
  \end{tabular}
  \end{center}
  よって $\sigma = (1, 2)$ が最小コスト $1$ を達成し，最適である．
\end{example}

\begin{claim}{最適解の存在}{sem-assignment-existence}
  任意の $n \in \N$ と任意の $\mathbf{C} \in \Rnn^{n \times n}$ に対して，
  最適割当問題（定義~\ref{def:sem-assignment}）の最小値は達成される．
\end{claim}

\begin{proof}
  $\Perm(n)$ は $|\Perm(n)| = n!$ の有限集合なので，集合
  \[
    \left\{ \frac{1}{n}\sum_{i=1}^{n} C_{i, \sigma(i)} \;\middle|\; \sigma \in \Perm(n) \right\}
  \]
  は高々 $n!$ 個の実数からなる空でない有限集合である．
  $\R$ の空でない有限部分集合は最小元を持つ（主張~\ref{clm:sem-finite-min}）
  ので，それを達成する $\sigma^* \in \Perm(n)$ が存在する．
\end{proof}

\begin{claim}{最適解が一意でない場合の存在}{sem-uniqueness}
  任意の $n \geq 2$ に対して，最適割当問題（定義~\ref{def:sem-assignment}）において
  異なる $\sigma_1, \sigma_2 \in \Perm(n)$ がともに最小値を達成するような
  $\mathbf{C} \in \Rnn^{n \times n}$ が存在する．
\end{claim}

\begin{proof}
  $C_{i,j} \defeq 1$（$i, j = 1, \ldots, n$）と定めると，任意の $\sigma \in \Perm(n)$ に対して
  $\frac{1}{n}\sum_{i=1}^{n} C_{i, \sigma(i)} = 1$．
  異なる $\sigma_1, \sigma_2 \in \Perm(n)$ を取れば，
  いずれも最小値 $1$ を達成する．
\end{proof}
